\documentclass[../../book]{subfiles}

\begin{document}
Error-correcting codes date back to 1948, when Shannon showed that information can be transmitted reliably over noisy channels, provided that redundancy is added~\cite{shannon1948}. A few years later, Hamming showed that one can build codes that not only detect possible errors but also correct them~\cite{hamming1950}.

The idea behind error-correcting codes is quite simple: encode a message of $k$ symbols into a longer \emph{codeword} of $n > k$ symbols, so that any two codewords differ in many positions. Then a few corrupted symbols cannot turn one codeword into another, and the original message can still be recovered.

The same idea turns out to be useful in cryptography. Shamir's secret sharing (\Cref{fig:shamir}) is, in fact, a Reed--Solomon code~\cite{mceliece1981}, and the Reed--Solomon fingerprinting from \Cref{section:toolkit} relies on the same property of polynomials. McEliece used codes to build one of the first public-key encryption schemes~\cite{McEliece78}, and code-based schemes are today among the main candidates for post-quantum cryptography.

In proof systems, codes play a different role. The simplest commitment that relies only on hash functions is a Merkle tree (\Cref{section:commitment-schemes}), but it gives the verifier only local access: it can open individual positions of the committed vector, and nothing more. If the prover committed to the message itself, a change in a single symbol would affect a single position, which the verifier would almost certainly miss. Committing to an encoding of the message with a code of large distance makes local access sufficient: a committed vector is either close to a single codeword, which binds the prover to a single message, or far from all codewords, which can be detected by checking only a few random positions using a test of proximity to the code. Many modern proof systems are built on this principle, including STARKs~\cite{EPRINT:BBHR18}, Ligero~\cite{CCS:AHIV17}, Brakedown~\cite{C:GLSTW23} and Binius~\cite{EC:DiaPos25}. Since they rely only on hash functions, they need no trusted setup and are believed to be secure against quantum adversaries.

In this section, we introduce linear codes and Reed-Solomon codes~\cite{guruswami23}, together with the notion of proximity to a code, which is the central tool behind STARKs (\Cref{section:stark}).

%This is precisely why we need to reformulate the problem: if we extend a polynomial over a domain much larger than its degree, any function that is not a valid polynomial will end up being 'far away' from actual polynomials. To leverage this property and solve our problem, we will turn to coding theory~\cite{guruswami23}.
%--------------------------------
\subsection{Linear Codes}
%--------------------------------
Let us first explain the basic idea of error-correcting codes.
\begin{intuition}
   Fix an alphabet $\Sigma$. Our goal is to send an encoded message $m \in \mathcal{A}^k$ over a noisy channel, which may corrupt some of the symbols. First, we use an encoder (encoding function) $\mathcal{E}$ to map $k$ digits to the $n$ encoded symbols, called a \emph{codeword}. In this way, we add $n-k$ redundant symbols to the message. Then, we transmitted the codeword over the noisy channel. The receiver gets a \emph{received word} of $n$ symbols, which may differ from the original codeword in some positions. The receiver then uses a decoder $\mathcal{D}$ to recover the original message $m$ from the received word. 
\end{intuition}

Now we can formalize the notion of a code and its encoding function.
\begin{definition}[Code]\label{def:code}
  A \emph{code} of block length $n$ over an alphabet $\Sigma$ is a subset $C \subseteq \Sigma^n$, whose elements are called \emph{codewords}. We write $M := |C|$ for the number of codewords and $k := \log_{|\Sigma|} M$ for the number of message symbols. An \emph{encoding} for $C$ is an injective map $E : \Sigma^k \to \Sigma^n$ whose image is $C$.
\end{definition}

Let us define a notion of distance that captures how close two vectors $\mathbf{u}$ and $\mathbf{v}$ are.
\begin{definition}[Hamming distance]\label{def:hamming-distance}
  Given $\mathbf{u}, \mathbf{v} \in \Sigma^n$, the \emph{Hamming distance} between $\mathbf{u}$ and $\mathbf{v}$, denoted as $\Delta(\mathbf{u}, \mathbf{v})$, is defined to be the number of positions in which they differ.
\end{definition}

To more clearly understand the Hamming distance, suppose we have two vectors $\mathbf{u}=00000$ and $\mathbf{v}=10100$. Then, the Hamming distance between $\mathbf{u}$ and $\mathbf{v}$ is $\Delta(\mathbf{u}, \mathbf{v})=2$ since they differ in the 1st and 3rd positions.

Hamming distance alone does not give you the full picture. For example, two vectors of length $1000$ that differ in $5$ places are much more similar than two vectors of length $10$ that also differ in $5$ places. Because of this, it is standard practice to normalize the distance by the block length.

\begin{definition}[Relative Hamming distance]\label{def:rel-distance}
  Let $\mathbf{u}, \mathbf{v} \in \Sigma^n$. The \emph{relative Hamming distance} between $\mathbf{u}$ and $\mathbf{v}$, denoted as $\delta(\mathbf{u}, \mathbf{v})$, is defined to be the Hamming distance between $\mathbf{u}$ and $\mathbf{v}$ divided by $n$:
  \begin{equation*}
    \delta(\mathbf{u}, \mathbf{v}) = \frac{1}{n} \Delta(\mathbf{u}, \mathbf{v}) \in [0,1]
  \end{equation*}
\end{definition}

While comparing individual vectors is helpful, a code's effectiveness depends on the entire set. Specifically, we care about the distance between codewords, where the worst-case spacing between any two codewords is the defining metric.
\begin{definition}[Minimum distance]\label{def:min-distance}
  Let $C\subseteq \Sigma^n$ be a code. The \emph{minimum distance} of $C$, denoted by $d_{\min}(C)$, is defined as 
  \begin{equation*}
    d_{\min}(C) = \min_{\mathbf{c}_1\neq \mathbf{c}_2 \in C} \Delta(\mathbf{c}_1, \mathbf{c}_2)
  \end{equation*}
    and the \emph{relative minimum distance} of $C$ is $\mu(C) := \frac{1}{n} d_{\min}(C)$.
\end{definition}

The minimum distance is the foundation of error correction. If $d_{\min}(C) = d$, introducing fewer than $d$ errors will never result in another valid codeword. Furthermore, with fewer than $d/2$ errors, the corrupted string remains strictly closer to the original codeword than to any other. Essentially, a larger minimum distance forces codewords further apart, so the original codeword can be recovered as long as fewer than $d/2$ errors occur.

\begin{example}\label{ex:repetition}
  The \emph{Repetition code} over $\mathbb{F}_2$ of block length $3$ repeats a single bit three times: $C=\{000, 111\}$. Here $k=1$ and $d_{\min}(C)=3$. If one bit of a codeword is flipped, e.g. $000 \mapsto 100$, the received word is still closer to $000$ than to $111$, so the error can be corrected since $d/2 = 3/2>1$. The opposite: if two bits are flipped, e.g. $000 \mapsto 110$, the received word is closer to $111$, and decoding fails since two errors are not fewer than $d/2$.
\end{example}

While calculating $d_{\min}(C)$ directly from the definition requires comparing every possible pair, practical codes offer a much faster shortcut. To introduce it, we first need to define the Hamming weight.

\begin{definition}[Hamming weight]\label{def:hamming-weight}
  The \emph{Hamming weight} of a vector $\mathbf{c} \in \Sigma^n$, denoted by $w(\mathbf{c})$, is defined as the number of non-zero symbols in $\mathbf{c}$.
\end{definition}

When $\Sigma$ is a field, these two concepts are directly linked by the equation $\Delta(\mathbf{u}, \mathbf{v}) = w(\mathbf{u} - \mathbf{v})$, which means that two vectors differ precisely in the coordinates where their difference is non-zero. 

In practice, cryptographic codes are never just arbitrary subsets of $\Sigma^n$ --- they are built with underlying algebraic structure. Assuming a code is closed under addition and scalar multiplication imposes virtually no practical limitations, yet it provides a twofold advantage.

\begin{definition}[Linear code]\label{def:linear-code}
  Let $\mathbb{F}_q$ be a finite field. A code $C \subseteq \mathbb{F}_q^n$ is a \emph{linear code} if it is a linear subspace of $\mathbb{F}_q^n$. If $C$ has dimension $k$ and minimum distance $d$, we denote it as an $[n,k,d]_q$ code.
\end{definition}

The first payoff is the shortcut promised above. Since the difference of two codewords of a linear code is again a codeword, the minimisation over all pairs collapses into a minimisation over single codewords.

\begin{proposition}\label{prop:dist-eq-weight}
  Let $C \subseteq \mathbb{F}_q^n$ be an $[n,k,d]_q$ code. Then
  \begin{equation*}
    d = \min_{\mathbf{0} \neq \mathbf{c} \in C} w(\mathbf{c}).
  \end{equation*}
\end{proposition}

\begin{proof}
  First, we show that $d$ is at most the minimum weight. Let $\mathbf{c}'$ be a non-zero codeword of minimum weight. Since $C$ is a linear subspace, $\mathbf{0} \in C$, and the distance from $\mathbf{c}'$ to $\mathbf{0}$ equals its weight. Hence $d \leq \Delta(\mathbf{0}, \mathbf{c}') = w(\mathbf{c}')$.

  Next, we show that $d$ is at least the minimum weight. Let $\mathbf{c}_1 \neq \mathbf{c}_2 \in C$ be such that $\Delta(\mathbf{c}_1, \mathbf{c}_2) = d$. By linearity, $\mathbf{c}_1 - \mathbf{c}_2 \in C$, and it is non-zero since $\mathbf{c}_1 \neq \mathbf{c}_2$. The non-zero symbols of $\mathbf{c}_1 - \mathbf{c}_2$ occur exactly in the positions where the two codewords differ, so $w(\mathbf{c}_1 - \mathbf{c}_2) = \Delta(\mathbf{c}_1, \mathbf{c}_2) = d$. Therefore the minimum weight of a non-zero codeword is at most $d$.
\end{proof}

The second payoff is that the dimension $k$ of a linear code is exactly the number of field elements each codeword carries, since $|C| = q^k$. This lets us put a price on the redundancy.

\begin{definition}[Rate]\label{def:rate}
  Let $C\subseteq \Sigma^n$ be a code. The \emph{rate} of $C$ is defined as 
  \begin{equation*}
    \rho = \frac{k}{n}\in [0,1]
  \end{equation*}
\end{definition}

A codeword consists of $n$ symbols but carries only $k$ symbols of information; thus $\rho$ is the fraction of each codeword that is actual payload. The rate and the relative minimum distance pull in opposite directions: adding redundancy lowers the rate but makes room for codewords that lie further apart.

\begin{example}[Parity code]\label{ex:parity}
  The \emph{parity code} $C \subseteq \mathbb{F}_2^n$ consists of all vectors with an even number of ones. It is a linear subspace of dimension $n-1$, since the last bit is determined by the others. Every non-zero codeword has weight at least $2$, and $110\ldots0 \in C$ has weight exactly $2$, so by \Cref{prop:dist-eq-weight} we get $d_{\min}(C) = 2$. Thus $C$ is an $[n, n-1, 2]_2$ code: its rate $\rho = (n-1)/n$ is close to $1$, but it can only \emph{detect} a single error, not correct it.
\end{example}

\begin{proposition}[Singleton bound]\label{prop:singleton}
  Every $[n,k,d]_q$ code satisfies $k \leq n-d+1$, equivalently $|C| \leq q^{n-d+1}$.
\end{proposition}

Codes attaining this bound are optimal and belong to a special class known as maximum distance separable.

\begin{definition}[MDS code]\label{def:mds}
  An $[n,k,d]_q$ code is called \textbf{maximum distance separable} (MDS) if $d = n-k+1$, that is, if it attains the Singleton bound.
\end{definition}

\begin{example}
  Both codes above are MDS: the repetition code of length $n$ is an $[n,1,n]_2$ code and the parity code is an $[n,n-1,2]_2$ code, and in both cases $d = n-k+1$. On the other hand, a $[7,4,3]_2$ code is not MDS, since $n-k+1 = 4 > 3$.
\end{example}
%--------------------------------
\subsection{Reed-Solomon Code}
%--------------------------------
We now introduce the primary code that forms the foundation for the rest of this section. Notably, it attains the Singleton bound.
\begin{definition}[Reed-Solomon Code]\label{def:rs}
  Let $\mathbb{F}_q$ be a finite field. Let $\bm{\alpha} = (\alpha_0, \alpha_1, \ldots, \alpha_{n-1})$ be a sequence of $n$ distinct elements (also called evaluation points) of $\mathbb{F}_q$ and choose $n$ and $k$ such that $k\leq n \leq q$. \emph{The Reed-Solomon} code is an encoding function $\mathsf{RS}_q[\bm{\alpha},k]: \mathbb{F}_q^k \to \mathbb{F}_q^n$ defined as follows: a message $\mathbf{m} = (m_0, m_1, \ldots, m_{k-1}) \in \mathbb{F}_q^k$ is mapped to a polynomial $\mathbf{m}\rightarrow f_{\mathbf{m}}(T)$, where
  \begin{equation*}
    f_{\mathbf{m}}(T) = \sum_{i=0}^{k-1} m_i\cdot T^i
  \end{equation*}
  Note that $f_{\mathbf{m}}(T)$ is a polynomial of degree at most $k-1$. The encoding of $\mathbf{m}$ is the evaluation of $f_{\mathbf{m}}(T)$ at all the $\alpha_i$'s:
\begin{equation*}
    \mathsf{RS}_q[\bm{\alpha},k](\mathbf{m}) = \big(f_{\mathbf{m}}(\alpha_0), f_{\mathbf{m}}(\alpha_1), \ldots, f_{\mathbf{m}}(\alpha_{n-1})\big).
\end{equation*}
\end{definition}

When $\bm{\alpha}$ and $k$ are clear from context, we simply write $\mathsf{RS}$.

\Cref{tab:rs-example} lists all codewords of the $[3,2,2]_3$ Reed--Solomon code whose evaluation points are $\mathbb{F}_3 = \{0,1,2\}$. The second row shows the polynomial $f_{\mathbf{m}}(T)$ corresponding to each message $\mathbf{m} \in \mathbb{F}_3^2$, and the codewords are ordered by the coefficients of these polynomials.

\begin{table}[h!]
\centering
\scriptsize
\begin{tabular}{|l|c|c|c|c|c|c|c|c|c|}
\hline
$\mathsf{RS}(\mathbf{m})$ & (0,0,0) & (1,1,1) & (2,2,2) & (0,1,2) & (1,2,0) & (2,0,1) & (0,2,1) & (1,0,2) & (2,1,0) \\ \hline
$f_{\mathbf{m}}(T)$       & $0$ & $1$ & $2$ & $T$ & $T+1$ & $T+2$ & $2T$ & $2T+1$ & $2T+2$ \\ \hline
\end{tabular}
\caption{All codewords of the $[3,2,2]_3$ Reed-Solomon code with evaluation points $\mathbb{F}_3$, together with the corresponding message polynomials.}
\label{tab:rs-example}
\end{table}

\begin{proposition}\label{prop:rs-linear}
  The Reed-Solomon code $\mathsf{RS}$ is a linear code.
\end{proposition}

\begin{proof}
    Let $a \in \mathbb{F}_q$ and $\mathbf{m}_1, \mathbf{m}_2 \in \mathbb{F}_q^k$. The polynomials $f_{\mathbf{m}_1}(T)$ and $f_{\mathbf{m}_2}(T)$ have degree at most $k-1$, and so do $a f_{\mathbf{m}_1}(T)$ and $f_{\mathbf{m}_1}(T) + f_{\mathbf{m}_2}(T)$. Since coefficients are added coordinate-wise:
    \begin{equation*}
      f_{\mathbf{m}_1}(T) + f_{\mathbf{m}_2}(T) = f_{\mathbf{m}_1+\mathbf{m}_2}(T),
    \end{equation*}
    and 
    \begin{equation*}
      a f_{\mathbf{m}_1}(T) = f_{a\mathbf{m}_1}(T)
    \end{equation*}

    In other words, the property of linearity holds for Reed-Solomon codes:
    \begin{gather*}
      \mathsf{RS}(\mathbf{m}_1) + \mathsf{RS}(\mathbf{m}_2) = \mathsf{RS}(\mathbf{m}_1+\mathbf{m}_2), \\ a \mathsf{RS}(\mathbf{m}_1) = \mathsf{RS}(a\mathbf{m}_1)
    \end{gather*}
    Therefore $\mathsf{RS}$ is an $[n,k]_q$ linear code.
\end{proof}

\begin{proposition}\label{prop:rs-distance}
  The Reed-Solomon code $\mathsf{RS}$ is an $[n,k,n-k+1]_q$ code. That is, it attains the Singleton bound.
\end{proposition}
\begin{proof}
  Let us fix arbitrary distinct $\mathbf{m}_1, \mathbf{m}_2\in \mathbb{F}_q^k$. Note that $f_{\mathbf{m}_1}(T), f_{\mathbf{m}_2}(T)\in \mathbb{F}_q[T]$ are distinct polynomials of degree at most $k-1$ since $\mathbf{m}_1\neq \mathbf{m}_2$. Then, the polynomial $f_{\mathbf{m}_1}(T)-f_{\mathbf{m}_2}(T)$ is also a non-zero polynomial of degree at most $k-1$. 

  Note that $w(\mathsf{RS}(\mathbf{m}_1)-\mathsf{RS}(\mathbf{m}_2))=\Delta(\mathsf{RS}(\mathbf{m}_1), \mathsf{RS}(\mathbf{m}_2))$. The weight of $\mathsf{RS}(\mathbf{m}_1)-\mathsf{RS}(\mathbf{m}_2)$ is $n$ minus the number of zeroes in $\mathsf{RS}(\mathbf{m}_1)-\mathsf{RS}(\mathbf{m}_2)$, which is equal to $n$ minus the number of roots that $f_{\mathbf{m}_1}(T)-f_{\mathbf{m}_2}(T)$ has in $\{\alpha_0, \alpha_1, \ldots, \alpha_{n-1}\}$. That is,
  \begin{equation*}
    \Delta(\mathsf{RS}(\mathbf{m}_1), \mathsf{RS}(\mathbf{m}_2)) = n - \#\{i \in [n] : f_{\mathbf{m}_1}(\alpha_i)=f_{\mathbf{m}_2}(\alpha_i)\}
  \end{equation*}
  Since $f_{\mathbf{m}_1}(T)-f_{\mathbf{m}_2}(T)$ is a non-zero polynomial of degree at most $k-1$, by \Cref{cor:root-count} it has at most $k-1$ roots in $\mathbb{F}_q$.
  Therefore, we have that the weight of $\mathsf{RS}(\mathbf{m}_1)-\mathsf{RS}(\mathbf{m}_2)$ is at least $n-(k-1)=n-k+1$. Therefore $d\geq n-k+1$, and since \Cref{prop:singleton} implies that $d\leq n-k+1$, we have that $d=n-k+1$.

  The argument above also shows that the distinct polynomials $f_{\mathbf{m}_1}(T)$ and $f_{\mathbf{m}_2}(T)$ are mapped to distinct codewords $\mathsf{RS}(\mathbf{m}_1)$ and $\mathsf{RS}(\mathbf{m}_2)$. Therefore, the code contains $q^k$ codewords, and has dimension $k$.
\end{proof}

\begin{remark}
  By \Cref{prop:rs-distance}, Reed-Solomon codes are MDS codes. In particular,
  $\mu(\mathsf{RS}) = 1 - \rho + \frac{1}{n}$, so lowering the rate directly increases the relative distance.
\end{remark}

%--------------------------------
\subsection{Proximity to a Code}\label{subsec:proximity}
%--------------------------------

The concepts discussed so far belong to classical coding theory, where the main objective is to recover the original message from a corrupted codeword. In our context, however, the requirement is weaker: the verifier does not actually need to know \emph{which} polynomial the prover committed to, but only whether the committed vector is close to \emph{some} valid codeword. To formalize this, we evaluate the distance between an arbitrary vector and the code as an entire set.

\begin{definition}[Distance to a code]\label{def:dist-to-code}
    Let $C \subseteq \Sigma^n$ be a code and $\mathbf{u} \in \Sigma^n$. The distance from
    $\mathbf{u}$ to $C$ is
    \begin{equation*}
        \Delta(\mathbf{u}, C) := \min_{\mathbf{c} \in C} \Delta(\mathbf{u}, \mathbf{c}),
    \end{equation*}
    and the relative distance is $\delta(\mathbf{u}, C) := \frac{1}{n}\Delta(\mathbf{u}, C)$.
\end{definition}

\begin{definition}[$\delta$-far]\label{def:delta-far}
    For $\delta \in [0,1]$, a vector $\mathbf{u} \in \Sigma^n$ is \textbf{$\delta$-far} from $C$
    if $\delta(\mathbf{u}, C) > \delta$, and \textbf{$\delta$-close} otherwise.
\end{definition}

\begin{example}
  Consider $\mathsf{RS}_5[\bm{\alpha},2]$ with $\bm{\alpha} = (0,1,2,3,4)$, that is, the evaluations of all polynomials of degree at most $1$ over $\mathbb{F}_5$. Let $\mathbf{u} = (0,1,4,4,1)$ be the evaluations of $T^2$. For every polynomial $g$ of degree at most $1$, the difference $T^2 - g(T)$ is a non-zero polynomial of degree $2$, so by \Cref{cor:root-count} it has at most $2$ roots. Hence $\mathbf{u}$ agrees with every codeword in at most $2$ positions, and $\Delta(\mathbf{u}, \mathsf{RS}) \geq 3$. The bound is attained by the constant codeword $(1,1,1,1,1)$, so $\delta(\mathbf{u}, \mathsf{RS}) = 3/5$, and $\mathbf{u}$ is, for instance, $1/2$-far from the code.
\end{example}

Why is this notion useful? Assume $\mathbf{u}$ is $\delta$-far from $C$. This implies that regardless of which codeword $\mathbf{c}$ we compare it to, the two will differ in strictly more than $\delta n$ coordinates. Consequently, a single uniformly random query will detect a discrepancy with a probability greater than $\delta$. If we make $t$ independent queries, the probability of missing the error entirely drops below $(1-\delta)^t$. Thus, a constant relative distance ensures that the soundness error decays exponentially with the number of queries. The fundamental challenge, however, is that the verifier does not know in advance which codeword $\mathbf{c}$ they are checking against. Constructing an efficient test that overcomes this limitation using only a small number of queries to $\mathbf{u}$ is exactly what the FRI protocol achieves.

%--------------------------------
\subsection{Reed-Solomon Proximity Problem}\label{subsec:rs-proximity}
%--------------------------------

We now state the problem that FRI solves~\cite{bensasson18}. Let $D \subseteq \mathbb{F}_q$ be a set of $n$ evaluation points, and let $\mathsf{RS}_q[D,k]$ denote the Reed-Solomon code whose evaluation points are the elements of $D$ (\Cref{def:rs}). The verifier has \emph{oracle access} to a function $f : D \to \mathbb{F}_q$: it cannot read $f$ as a whole, but it may query its value at any point of $D$.
In practice, the prover commits to $f$ with a Merkle tree (\Cref{section:commitment-schemes}), so that each query costs one opening.

\begin{definition}[Reed-Solomon proximity problem]\label{def:rs-proximity}
  Given oracle access to $f : D \to \mathbb{F}_q$ and $\delta \in (0,1)$, distinguish, with high confidence and few queries, between the case $f \in \mathsf{RS}_q[D,k]$ and the case that $f$ is $\delta$-far from $\mathsf{RS}_q[D,k]$.
\end{definition}

When the verifier is restricted to oracle access to $f$, the problem is known as \emph{Reed-Solomon proximity testing}. In this model, exactly $k+1$ queries are both necessary and sufficient: the verifier reads $k$ points to interpolate the unique polynomial of degree less than $k$, and then checks it for consistency against one additional random point. However, since $k$ is typically extremely large in real-world applications, this straightforward interpolation is computationally infeasible.

To overcome this limitation, we allow the prover to participate. In \emph{Reed-Solomon proximity verification}, the prover and verifier interact via an IOP (\Cref{fig:iop}). During each round, the prover commits to an additional oracle, and the verifier is allowed to query any of them. This interactive setup is known as an \emph{IOP of proximity} (IOPP), making standard proximity testing just the degenerate case where the prover provides no auxiliary information.

The \emph{Fast Reed-Solomon IOPP} (FRI) is precisely this type of protocol. It achieves linear prover complexity in $n$, while keeping both the verifier's computational effort and the query complexity strictly logarithmic in $n$.

\subsection{Exercises}

\begin{xexercise}
  {Exercise 1.}
  {Over $\mathbb{F}_3$, consider $\mathbf{u} = (1,0,2,2,1)$ and $\mathbf{v} = (1,2,2,0,1)$. What are the
  Hamming distance $\Delta(\mathbf{u},\mathbf{v})$ and the relative Hamming distance $\delta(\mathbf{u},\mathbf{v})$?}
  {2}
  {
    \item $\Delta = 2$, $\delta = 2/5$
    \item $\Delta = 3$, $\delta = 3/5$
    \item $\Delta = 2$, $\delta = 2/3$
    \item $\Delta = 4$, $\delta = 4/5$
  }
\end{xexercise}

\begin{xexercise}
  {Exercise 2.}
  {Let $C \subseteq \mathbb{F}_2^4$ be the linear code spanned by $(1,1,0,1)$ and $(0,1,1,1)$. What is its
  minimum distance $d_{\min}(C)$?}
  {4}
  {
    \item $1$
    \item $2$
    \item $3$
    \item $4$
  }
\end{xexercise}

\begin{xexercise}
  {Exercise 3.}
  {Which of the following parameters $[n,k,d]_q$ cannot be achieved by any linear code?}
  {4}
  {
    \item $[7,4,4]_q$
    \item $[7,4,5]_q$
    \item $[5,1,5]_q$
    \item $[6,3,4]_q$
  }
\end{xexercise}

\begin{xexercise}
  {Exercise 4.}
  {Consider the Reed--Solomon code $\mathsf{RS}_7[\bm{\alpha},2]$ with $\bm{\alpha} = (0,1,2,3,4,5,6)$.
  What is the encoding of the message $\mathbf{m} = (3,2)$?}
  {2}
  {
    \item $(2,5,1,4,0,3,6)$
    \item $(3,5,7,9,11,13,15)$
    \item $(3,5,0,2,4,6,1)$
    \item $(3,2,0,0,0,0,0)$
  }
\end{xexercise}

\begin{xexercise}
  {Exercise 5.}
  {Consider $\mathsf{RS}_7[\bm{\alpha},2]$ with $\bm{\alpha} = (0,1,\ldots,6)$, and let
  $\mathbf{u} = (0,1,4,2,2,4,1)$ be the evaluations of $T^2$. What is $\delta(\mathbf{u}, \mathsf{RS})$?}
  {4}
  {
    \item $2/7$
    \item $3/7$
    \item $5/7$
    \item $6/7$
  }
\end{xexercise}

\begin{xexercise}
  {Exercise 6.}
  {Suppose $\mathbf{u}$ is $1/2$-far from a code $C$, and the verifier compares $\mathbf{u}$ with a fixed
  codeword at $t$ independent uniformly random positions. What is the smallest $t$ for which the bound
  $(1-\delta)^t$ guarantees that the probability of missing every discrepancy is at most $2^{-40}$?}
  {4}
  {
    \item $20$
    \item $40$
    \item $80$
    \item $128$
  }
\end{xexercise}

\end{document}